% STATUS: written 2026-10-21, recalled results as is (one environment each, source kind and whole statement)
% Sources read in full: Herstein 1957 (pp. 1104--1110), Cusack 1975 (pp. 321--324), Bresar and Vukman 1988 (pp. 321--322).
\section{The theorem of Herstein and its proofs}\label{sec:jordan-type:herstein-theorem}

Chapter~\ref{ch:prime-semiprime} ended with Jordan ideals of a prime ring. We now
ask what the Jordan product $a\circ b=ab+ba$ does to the product rule itself.

Herstein begins with this product. From an associative ring $A$ one builds the
Jordan ring of $A$ by defining $a\circ b=ab+ba$ \cite[p.~1104]{Herstein1957}.
A \emph{Jordan derivation} of $A$ is a derivation of this Jordan ring. It is an additive map $'$
with $(a\circ b)'=a'\circ b+a\circ b'$ \cite[p.~1104]{Herstein1957}. Written out, this
is $(ab+ba)'=a'b+ab'+b'a+ba'$ \cite[p.~1104]{Herstein1957}. Every derivation is a
Jordan derivation. Herstein asks about the converse. He records that
Kaplansky first raised the idea of a Jordan derivation in conversation
\cite[p.~1104, footnote]{Herstein1957}.

The obvious rival is a map $*$ with $(ab)^*=b^*a+ba^*$, which Herstein calls a
\emph{reverse derivation} \cite[pp.~1104--1105]{Herstein1957}. Such a map would
be a Jordan derivation. Herstein first shows that it hardly ever occurs.

\begin{theorem}[{\cite[Theorem~2.1, p.~1105]{Herstein1957}}]\label{thm:jordan:herstein-reverse}
Let $A$ be a prime ring and suppose that $*$ is a nonzero reverse derivation of $A$.
Then $A$ is a commutative integral domain and $*$ is an ordinary derivation of $A$.
\end{theorem}

The main result needs two lemmas. In Section~3 of his paper Herstein assumes that $A$
is a prime ring of characteristic different from $2$.

\begin{lemma}[{\cite[Lemma~3.1, p.~1106]{Herstein1957}}]\label{lem:jordan:herstein-aba}
If $'$ is a Jordan derivation of $A$, then for all $a,b\in A$,
$(aba)'=a'ba+ab'a+aba'$.
\end{lemma}

The next lemma records the key identity of the proof.

\begin{lemma}[{\cite[Lemma~3.3, p.~1107]{Herstein1957}}]\label{lem:jordan:herstein-key}
For all $a,b\in A$, $((ab)'-a'b-ab')(ab-ba)=0$.
\end{lemma}

Herstein proves Lemma~\ref{lem:jordan:herstein-key} by expanding $(ab(ab)+(ab)ba)'$ in
two ways \cite[p.~1107]{Herstein1957}. He then treats commuting and noncommuting
pairs $a,b$ separately \cite[pp.~1107--1110]{Herstein1957}. The fact that $A$ is prime
finishes the argument.

\begin{theorem}[{\cite[Theorem~3.1, p.~1110]{Herstein1957}}]\label{thm:jordan:herstein}
If $A$ is a prime ring of characteristic different from $2$, then any Jordan
derivation of $A$ is an ordinary derivation of $A$.
\end{theorem}

Herstein notes that the characteristic is used in only two places: in the proof of
Lemma~\ref{lem:jordan:herstein-aba} and in the commutative case
\cite[p.~1110]{Herstein1957}. So he changes the definition of a Jordan derivation.
He asks that $(a+b)'=a'+b'$, $(a^{2})'=aa'+a'a$ and $(aba)'=a'ba+ab'a+aba'$. For
characteristic different from $2$ the third condition follows from the first two
\cite[p.~1110]{Herstein1957}. With this definition he treats characteristic~$2$.

\begin{theorem}[{\cite[Theorem~4.1, p.~1110]{Herstein1957}}]\label{thm:jordan:herstein-char-two}
Let $A$ be a prime ring of characteristic $2$. Then, if $A$ is not a commutative integral
domain, any Jordan derivation of $A$ is an ordinary derivation.
\end{theorem}

Eighteen years later Cusack dropped the hypothesis that the ring is prime. He works
in a ring $R$ in which $2x=0$ implies $x=0$, and he wants to enlarge the class of
rings for which the converse holds \cite[p.~321]{Cusack1975}. For a Jordan derivation
$D$ he sets $d(a,b)=D(ab)-aDb-(Da)b$ \cite[p.~321]{Cusack1975}. His main step is that
$[a,b]\circ d(a,b)$ and $[[a,b],r,d(a,b)]$ vanish \cite[Lemma~2, p.~322]{Cusack1975}.
He gets this from the same identities for Jordan homomorphisms
\cite[Lemma~1, p.~322]{Cusack1975}.

\begin{theorem}[{\cite[Theorem~4, pp.~322--323]{Cusack1975}}]\label{thm:jordan:cusack-radical}
Let $D$ be a Jordan derivation on a ring $R$ in which $2x=0$ implies $x=0$. Then for
all $a$ and $b$ in $R$, $D(ab)-aDb-(Da)b$ is in $L$, the Baer lower radical of $R$.
\end{theorem}

A ring is semiprime when its lower radical is zero \cite[p.~322]{Cusack1975}. Cusack
states two consequences, the first in his setting.

\begin{corollary}[{\cite[Corollary~5, p.~323]{Cusack1975}}]\label{cor:jordan:cusack-semiprime}
If $R$ is semiprime, then $D$ is a derivation.
\end{corollary}

The next theorem states the second consequence, for a ring with a commutator that is not a zero divisor.

\begin{theorem}[{\cite[Theorem~6, p.~323]{Cusack1975}}]\label{thm:jordan:cusack-commutator}
Let $R$ be a ring where $2x=0$ implies $x=0$ and which has a commutator which is not a
zero divisor. Then every Jordan derivation on $R$ is a derivation.
\end{theorem}

Cusack also shows that some hypothesis is needed. Take a ring with an element $x$ such
that $axa=0$ for all $a$ but $axb\neq0$ for some $a,b$. Then $Da=ax$ is a Jordan
derivation that is not a derivation \cite[\S3, p.~324]{Cusack1975}. Now take a ring with
$a^{2}=0$ for all $a$ and $ab\neq0$ for some $a,b$. The Jordan product is then zero, so
every additive map is a Jordan derivation, and the identity map is not a derivation
\cite[\S3, p.~324]{Cusack1975}. He leaves open the search for necessary conditions
\cite[p.~324]{Cusack1975}.

Finally, a short proof of Theorem~\ref{thm:jordan:herstein} appeared in 1988. Brešar and
Vukman write $a^{b}$ for $D(ab)-D(a)b-aD(b)$ \cite[p.~321]{BresarVukman1988}. For a
Jordan derivation $D$, they note that $b^{a}=-a^{b}$ \cite[p.~321]{BresarVukman1988}.
Their key identity follows.

\begin{theorem}[{\cite[Theorem~3, p.~322]{BresarVukman1988}}]\label{thm:jordan:bresar-vukman}
Let $R$ be a ring of characteristic not two, and let $D\colon R\to R$ be a Jordan
derivation. In this case for all $a,b,r\in R$ we have
$a^{b}r(ab-ba)+(ab-ba)ra^{b}=0$.
\end{theorem}

They prove it by expanding $abrba+barab$ in two ways \cite[p.~322]{BresarVukman1988}.
With Lemma~3.10 of the \emph{Topics in ring theory} of Herstein they obtain
Theorem~\ref{thm:jordan:herstein} \cite[p.~322]{BresarVukman1988}.

Section~\ref{sec:jordan-type:semiprime} follows the semiprime case further.
